
\begin{proposition}[Every nonlinear primary Hopf surface occurs unchanged over a rational base]\label{cdp:st:C.2.UN.42}
For every integer $r\ge1$ and every $\beta$ with $0<|\beta|<1$,
the primary Hopf surface
\[
 H_r(\beta)=(\C^2\setminus\{0\})/
   \langle(z,w)\mapsto(\beta^r z+w^r,\beta w)\rangle
                                                        \tag{C.NH.1}\label{cdp:eq:C.NH.1}
\]
occurs unchanged as a smooth multiplicity-one component of a
reduced SNC fibre of a proper map $f:Y\to\PP^1$, where $Y$ is
a smooth compact complex threefold and the general fibre is a
complex two-torus. Every reduced fibre support in this
construction has Euler number zero.

The distinguished full fibre is $S=H_r(\beta)+P_r$, with one
smooth elliptic intersection $J$. The second component is a
smooth elliptic surface over an elliptic curve, with exactly
one multiple fibre $rJ$ when $r>1$, and
\[
 \begin{gathered}
 K_{P_r}=\OO_{P_r}((r-1)J),\qquad
 b(P_r)=(1,3,4,3,1),\qquad h^\bullet(\OO_{P_r})=(1,2,1),\\
 H^\bullet(S,\Z)\text{ is torsion-free of ranks }(1,3,5,5,2),
 \qquad h^\bullet(S,\OO_S)=(1,2,1).
 \end{gathered}                                         \tag{C.NH.2}\label{cdp:eq:C.NH.2}
\]
The Hopf component has integral ranks $(1,1,0,1,1)$ and
coherent dimensions $(1,1,0)$. The construction works for
every prescribed $\beta$, without a smallness condition
beyond $|\beta|<1$.
\end{proposition}

\begin{proof}
\par\medskip\noindent\textbf{A compact base with the required translation jet.}
Let $B$ be the elliptic cubic $v^2=u^3-u+1$, with origin the
point at infinity. Its cubic has distinct nonzero roots. Let
$C$ be the smooth projective completion of
\[
 v^2=t^{6r}-t^{2r}+1,
 \qquad \psi:C\longrightarrow B,\quad
          (t,v)\longmapsto(t^{2r},v).
                                                        \tag{C.NH.3}\label{cdp:eq:C.NH.3}
\]
The polynomial has distinct roots: a repeated root would be
nonzero and would give a common root of $u^3-u+1$ and its
derivative. At infinity the two charts
$s=t^{-1}$, $v s^{3r}=\pm\sqrt{1-s^{4r}+s^{6r}}$ give two
distinct smooth points $c_+,c_-$. Thus $C$ is a smooth compact
curve of genus $3r-1$. The involution $h_C(t,v)=(t,-v)$ has
quotient $\PP^1_t$, interchanges $c_+,c_-$, and satisfies
$\psi h_C=-\psi$. Its fixed points are the $6r$ finite roots
of the displayed polynomial.

The elliptic uniformizer $u/v$ has order $r$ at both
$c_\pm$. Choose a local additive coordinate $p$ on $B$ at
zero. At each of the two points a holomorphic coordinate
$s$ on $C$ can therefore be chosen so that
\[
 2\psi(s)=s^r\quad\hbox{in the additive coordinate on }B.
                                                        \tag{C.NH.4}\label{cdp:eq:C.NH.4}
\]
Taking the $r$-th root of its nonzero holomorphic unit proves
this coordinate equality. On the compact surface $B\times C$
consider the commuting automorphisms
\[
 A(p,c)=(p+2\psi(c),c),\qquad
 A_{1/2}(p,c)=(p+\psi(c),c),\qquad
 h(p,c)=(-p,h_Cc).
                                                        \tag{C.NH.5}\label{cdp:eq:C.NH.5}
\]
Here $A_{1/2}^2=A$, $h^2=1$, and the equality
$\psi h_C=-\psi$ proves that $h$ commutes with both
translations. All three preserve the pair of points
$(0,c_+),(0,c_-)$.

\par\medskip\noindent\textbf{The exact weighted blowup and its smooth nonzero line.}
Perform the weighted blowup of weights $(r,1)$ at both
points, retaining its smooth orbifold charts. Denote the
resulting compact orbifold surface by $\mathcal V$, its
exceptional divisors by $R_+,R_-$, and put
$L=\OO_{\mathcal V}(R_++R_-)$. Locally its weighted blowup
has the explicit presentation
\[
 \left[((\C^2_{z,w}\setminus\{0\})\times\C_x)/\C^*\right],
 \quad \lambda(z,w,x)=(\lambda^r z,\lambda w,\lambda^{-1}x),
 \quad p=x^r z,\quad s=xw.
                                                        \tag{C.NH.6}\label{cdp:eq:C.NH.6}
\]
The exceptional divisor is $x=0$. The line $L$ has fibre
coordinate $\ell$ of weight $-1$ and canonical section
$\ell=x$. The quotient chart $w\ne0$ is the smooth chart
$p=s^r(z/w^r)$; the chart $z\ne0$ has the residual finite
action
\[
 (w,x,\ell)\longmapsto
      (\zeta w,\zeta^{-1}x,\zeta^{-1}\ell),\qquad \zeta^r=1.
                                                        \tag{C.NH.7}\label{cdp:eq:C.NH.7}
\]
It is free when $\ell\ne0$. More explicitly, on $L^\times$
use $\lambda=\ell$ in \eqref{cdp:eq:C.NH.6} to set $\ell=1$. This gives
the actual smooth coordinates
$(z,w,x)\in(\C^2\setminus\{0\})\times\C$, after replacing
the coordinates by $(\ell^rz,\ell w,\ell^{-1}x)$.
Thus the coarse nonzero line space $L^\times$ is a smooth
complex threefold, including the orbifold points.

All three maps in \eqref{cdp:eq:C.NH.5} lift canonically. In the
coordinates \eqref{cdp:eq:C.NH.4}, $A$ is $p\mapsto p+s^r$ and its lift
in \eqref{cdp:eq:C.NH.6} is $(z,w,x)\mapsto(z+w^r,w,x)$.
The weighted ideals are therefore preserved. The half
translation replaces the coefficient by $1/2$; $h$
interchanges the two centres and preserves their weighted
orders. The line lifts preserve its canonical divisor
section, so their group laws hold exactly, including
$\widetilde A_{1/2}^{\,2}=\widetilde A$ and
$\widetilde h^{\,2}=1$.

Multiply the line fibre by $\beta$ and compose with
$\widetilde A$. On the smooth coordinates just constructed
the resulting automorphism is
\[
 g(z,w,x)=
       (\beta^r(z+w^r),\beta w,\beta^{-1}x),
 \qquad s=xw.
                                                        \tag{C.NH.8}\label{cdp:eq:C.NH.8}
\]
It preserves the map to $C$. On $x=0$, the change
$Z=\beta^{-r}z$ turns its first coordinate into
$\beta^r Z+w^r$, proving the exact identification with
\eqref{cdp:eq:C.NH.1}.

\par\medskip\noindent\textbf{Uniform properness for every contraction.}
We verify the global quotient, rather than inferring its
compactness from a local contraction. Choose a Hermitian
orbifold metric on $L$. Write $\sigma$ for its canonical
section. Near either centre, its norm off the exceptional
divisor is comparable, with constants independent of the
point, to
\[
       \rho(p,s)=\bigl(d_B(p,0)^2+|s|^{2r}\bigr)^{1/(2r)}.
                                                        \tag{C.NH.9}\label{cdp:eq:C.NH.9}
\]
Here $d_B$ is a fixed flat elliptic distance. To see the
comparison, normalize the weighted coordinates so that
$|z|^2+|w|^{2r}=1$. Then
$(|p|^2+|s|^{2r})^{1/(2r)}=|x|$, which is the norm of the
canonical section up to a smooth positive bounded factor
on the compact weighted exceptional curve. Away from the
two centres all these norms are bounded above and below.
More precisely this comparison holds on $B$ times a small
$c_\pm$-disc: outside the chosen $p$-coordinate neighborhood,
both $\rho$ and $\|\sigma\|$ are bounded above and away from
zero. It therefore also covers an iterate $A^n$ that moves
$p$ out of that local coordinate chart.

For every integer $n$, the triangle inequality on the flat
elliptic curve gives
\[
 d_B(p+2n\psi(s),0)
       \le d_B(p,0)+C_0|n||s|^r.
\]
The same inequality applied to the inverse translation
gives the reverse comparison. Since the canonical lift
sends $\sigma(p,s)$ to $\sigma(p+2n\psi(s),s)$, there is
a global constant $C_1\ge1$ such that
\[
 C_1^{-1}(1+|n|)^{-1/r}\|\ell\|
 \le \|\widetilde A^{\,n}\ell\|
 \le C_1(1+|n|)^{1/r}\|\ell\|.
                                                        \tag{C.NH.10}\label{cdp:eq:C.NH.10}
\]
The comparisons first hold off the exceptional divisors.
The ratio of line norms extends continuously and positively
on the entire orbifold line, so they hold there as well.
On the compact part of $C$ away from both centres the
canonical section is nowhere zero and supplies the same
bounds with a fixed constant. This proves the uniform
global assertion.

Consequently $\|g^n\ell\|$ tends uniformly to zero on the
unit circle bundle as $n\to+\infty$, and to infinity as
$n\to-\infty$: multiply \eqref{cdp:eq:C.NH.10} by $|\beta|^n$.
The action is properly discontinuous because any two
compact subsets have line norms bounded above and away
from zero. It is free, since a fixed point of a nonzero
power would also be fixed by arbitrarily large positive
powers, contradicting the norm limit.

Choose $N$ such that $g^N$ decreases every norm by a factor
strictly less than one. Its minimum norm ratio is a positive
number $c_N$. Every $g^{N\Z}$-orbit meets the compact annulus
$c_N\le\|\ell\|\le1$: take the first iterate whose norm is
at most one. This annulus is compact over the compact
orbifold base. It follows that
\[
                 \widetilde Y=L^\times/\langle g\rangle
                                                        \tag{C.NH.11}\label{cdp:eq:C.NH.11}
\]
is a smooth compact threefold, with a proper map to $C$.
This proof applies to every $0<|\beta|<1$.

Off the two exceptional fibres, the canonical section
trivializes $L$ and gives the actual quotient
\[
 (B\times\C^*)/
   \langle(p,\ell)\mapsto(p+2\psi(c),\beta\ell)\rangle.
                                                        \tag{C.NH.12}\label{cdp:eq:C.NH.12}
\]
If $B=\C/\Lambda_B$, its additive covering lattice is
generated by $(\lambda,0)$ for $\lambda\in\Lambda_B$,
$(0,2\pi i)$, and $(a_c,\log\beta)$, where $a_c$ is any
lift of $2\psi(c)$. These four real generators are
independent because $\log|\beta|\ne0$. Changing the lift
of $a_c$ changes a generator by $\Lambda_B$. This proves
that \eqref{cdp:eq:C.NH.12} is precisely a compact two-torus.

At either distinguished value the exact function is
$s=xw$. Its full fibre is reduced SNC, with components
$H:x=0$ and $P:w=0$, meeting along
\[
          J=(\C^*_z/\beta^{r\Z}),\qquad x=w=0.
                                                        \tag{C.NH.13}\label{cdp:eq:C.NH.13}
\]
Both components and their intersection are compact since
they are closed in \eqref{cdp:eq:C.NH.11}. On $P$ the map to the
unchanged elliptic base $B$ has the exact local form
\[
  p=x^r z,\qquad (z,x)\sim(\beta^r z,\beta^{-1}x).
                                                        \tag{C.NH.14}\label{cdp:eq:C.NH.14}
\]
Away from $p=0$ its canonical frame is the product
$B\setminus\{0\}\times E_\beta$; on the overlap the
frame coordinate is $\ell=x^{-1}$ and $z=p\ell^r$.
These equalities identify the surface with the one-fibre
logarithmic model in (C.ER.10)--(C.ER.16), with $m=r$.
They also give the same model for $r=1$ directly.
In particular every fibre support of $P\to B$ is a
smooth elliptic curve, and the special full fibre is
$rJ$. Stratified Euler additivity gives $e(P)=0$.
The Hopf surface and $J$ also have Euler number zero,
so the entire distinguished support has Euler number zero.
The local holomorphic canonical evaluation form is
$x^r\,dz\wedge dw\wedge dx$. The Jacobian of \eqref{cdp:eq:C.NH.8}
is $\beta^r$, and the multiplier of $x^r$ is $\beta^{-r}$,
so this form is invariant. Its holomorphic vanishing orders
along $(H,P)$ are exactly $(r,0)$. This realizes the
strict canonical-order difference in \ref{cdp:st:C.2.UN.40} without
replacing a holomorphic generator by a meromorphic
full-fibre shift.

\par\medskip\noindent\textbf{The free quotient giving the rational base.}
Fix either square root $\beta^{1/2}$. The commuting
canonical lifts define
\[
           G=\beta^{1/2}\widetilde A_{1/2}\widetilde h,
                    \qquad G^2=g.
                                                        \tag{C.NH.15}\label{cdp:eq:C.NH.15}
\]
Thus $G$ induces an involution on $\widetilde Y$. A fixed
point there would have $h_Cc=c$. Such a point is not one
of the two centres. Moreover $\psi(c)$ is a two-torsion
point of $B$, so $2\psi(c)=0$. In the canonical line frame,
\[
 G(p,c,\ell)=(-p-\psi(c),c,\beta^{1/2}\ell),
 \qquad g^n(p,c,\ell)=(p,c,\beta^n\ell).
                                                        \tag{C.NH.16}\label{cdp:eq:C.NH.16}
\]
The equality $G(p,c,\ell)=g^n(p,c,\ell)$ is impossible by
the absolute values of the line coordinates: it would
give $|\beta|^{1/2}=|\beta|^n$. The involution is therefore
free, and
\[
                 Y=\widetilde Y/\langle G\rangle
                  \longrightarrow C/\langle h_C\rangle=\PP^1
                                                        \tag{C.NH.17}\label{cdp:eq:C.NH.17}
\]
is a smooth compact threefold with the required map.
The two Hopf components over $c_+,c_-$ are exchanged.
Restriction of the quotient to either one is an
isomorphism onto its image, so the prescribed Hopf
surface is unchanged. The entire distinguished fibre is
isomorphic to either $H+P$ above, and remains reduced.
At a fixed value of $h_C$ the full fibre is twice a free
quotient of a two-torus; its reduced support has Euler
number zero. At every other ordinary value it is a
two-torus. All fibres are connected, as is seen directly
from these quotient descriptions and from $H\cap P=J$.
This proves all the stated occurrence and Euler assertions.

\par\medskip\noindent\textbf{The component and support groups.}
The exact local frame in \eqref{cdp:eq:C.NH.14} gives a global canonical
form $dp\wedge d\log\ell$ off $J$. In that chart it is
$-x^{r-1}dz\wedge dx$, and it is invariant under the
displayed deck map. Hence its divisor is $(r-1)J$.
The direct-image equality proved from this same chart is
$p_*\OO_P(kJ)=\OO_B(\lfloor k/r\rfloor[0])$.
It gives $h^0(K_P)=1$. Since $J^2=0$ and $e(P)=0$,
surface Riemann--Roch gives $\chi(\OO_P)=0$ and
$h^\bullet(\OO_P)=(1,2,1)$.

For the integral groups use the exact Seifert filling in
(C.ER.16). Its meridian is $r\delta+h$ and its longitude
is $\delta$, where $\delta$ is the boundary of the
once-punctured elliptic base and $h$ the product circle.
The base relation kills $\delta$ in first homology, and
the meridian then kills $h$. Thus the closed oriented
three-manifold $M_r$ has $H_1(M_r,\Z)=\Z^2$; all its
homology is free, by integral Poincare duality and the
universal coefficient sequence. The radial line coordinate
gives $P\simeq M_r\times S^1$ smoothly. For a nonreal
$\beta$ its angular rotation is isotopic to the identity
through the same circle action, so the mapping torus has
this same product identification. This proves the ranks
of $P$ in \eqref{cdp:eq:C.NH.2}.
That circle isotopy preserves the Seifert core. Thus it
identifies the pair $(P,J)$ with
$(M_r\times S^1,\mathrm{core}\times S^1)$, not merely
the ambient cohomology groups.

The angular loop of $J$ is the solid-torus longitude and
is zero in $H_1(M_r,\Z)$; its radial loop maps primitively
to the last $S^1$. Therefore $[J]=0$ in $H_2(P,\Z)$,
including integrally. In the Hopf component its angular
loop is zero and its radial loop is the primitive deck
generator. The integral normalization sequence of
$S=P\cup_JH$ consequently has the restriction maps
\[
 H^0(P\sqcup H,\Z)\to H^0(J,\Z):\ (a,b)\mapsto a-b,
 \qquad
 H^1(P\sqcup H,\Z)\to H^1(J,\Z):
 \begin{pmatrix}0&0&1&-1\\0&0&0&0\end{pmatrix},
                                                        \tag{C.NH.18}\label{cdp:eq:C.NH.18}
\]
up to simultaneous orientation choices for the radial
loops, and its degree-two map is zero. Both displayed
images are primitive. The exact sequence gives precisely
the free groups of ranks $(1,3,5,5,2)$.

Finally the constant deck cocycle represents the nonzero
class in $H^1(H,\OO_H)$. It is not a coboundary: a
holomorphic solution of $a\circ F-a=1$ on $\C^2\setminus0$
extends across zero and gives $0=1$ there. Its restriction
to $J$ is still nonzero, since the constant Laurent
coefficient of $a(\beta^r z)-a(z)$ is zero. Thus the
coherent degree-one restriction onto $H^1(J,\OO_J)$
is surjective. Together with the degree-zero surjection
and the component groups already computed, the coherent
normalization sequence gives $h^\bullet(S,\OO_S)=(1,2,1)$.
All support maps are the actual inclusions in \eqref{cdp:eq:C.NH.13}--\eqref{cdp:eq:C.NH.14}.
\end{proof}


